\section{Literature Review and Foundations}

\subsection{Regulatory Principles and Technical Signals}

The GDPR regulates processing rather than permission counts. Article 5 establishes principles including lawfulness, transparency, purpose limitation, data minimisation, accuracy, storage limitation, integrity, confidentiality, and accountability. Articles 6 and 9 concern lawful bases and special-category data; Article 25 requires data protection by design and by default \cite{eurlex2016}. EDPB guidance describes Article 25 as a continuing obligation implemented through appropriate technical and organisational measures \cite{edpb2019}. A technical observation may support an accountability inquiry, but cannot by itself demonstrate whether these legal conditions are satisfied.

This distinction determines the framework's vocabulary. A threshold exceedance is a \emph{review signal}. Missing purpose, lawful-basis, necessity, data-subject, controller, retention, and sharing evidence is shown explicitly. The UI never derives a finding labelled ``GDPR violation'' or ``compliant''. Regulation-pack references provide a reason to investigate and a link to an official source; they do not replace legal interpretation.

\subsubsection{From Legal Principle to Review Question}

Operationalisation is necessary in any regulatory technology, but it must preserve the semantic distance between a principle and a sensor. Article 5(1)(a) concerns lawfulness, fairness, and transparency. A permission count can contribute to transparency by revealing a pattern that was previously invisible, yet it cannot establish lawfulness or fairness. Article 5(1)(b) concerns purpose limitation. A count contains no purpose unless purpose is supplied as contextual evidence. Article 5(1)(c) concerns data minimisation. Frequency can be relevant to necessity, but a high frequency may be required for continuous navigation while a single collection may be excessive for an unrelated purpose. Article 5(2) concerns accountability. An audit record can support demonstration, but it is only one component of the records, decisions, safeguards, and governance required from a controller \cite{eurlex2016}.

Articles 6 and 9 further show why a permission signal cannot become a verdict. Article 6 provides several possible lawful bases; consent is not automatically the only basis. Article 9 establishes a prohibition and exceptions for special categories, but not every Android dangerous permission maps directly to special-category processing. Location, contacts, microphone, and camera data are context-dependent. Audio might contain health information, ordinary conversation, environmental noise, or no retained personal data. A technically observed invocation does not identify which of those states occurred.

The pack therefore maps a signal to a sequence of questions. For a location-related pattern, the reviewer may ask what purpose was declared, whether continuous precision was necessary, whether a less intrusive source was available, whether the application was in the foreground, whether location was retained or transmitted, and whether the user could reasonably anticipate the flow. For a microphone-related pattern, the reviewer may ask whether recording was user-initiated, how activation was signalled, whether raw audio left the device, and how long any derivative was retained. These questions are not answers; their value is that they prevent the visible count from crowding out legally decisive facts.

\subsubsection{Data Protection by Design and Interface Conduct}

Article 25 directs controllers to implement appropriate measures and privacy-protective defaults, taking account of state of the art, cost, context, and risk. EDPB guidance emphasises that this is a continuing obligation rather than a certification achieved once \cite{edpb2019}. For Privacy Lens, data protection by design affects both internal processing and the claims made to users. Local storage, disabled backup, no analytics, no advertising, no account, explicit deletion, and minimised permissions reduce the application's own data footprint. Fail-closed validation and immutable pack metadata reduce the risk that incomplete evidence becomes a misleading result.

Interface conduct is also relevant. A privacy tool could employ alarming colours, default users into sharing diagnostic data, hide evidence limitations, or imply that purchasing a service resolves risk. EDPB guidance on deceptive design patterns provides a broader reminder that formal information can still be presented in a manipulative way \cite{edpb2022}. Privacy Lens uses restrained severity language, places limitations near the result, and keeps the demonstration optional. The implementation is not claimed to have been audited against every deceptive-design criterion, but the guidance informs the design objective of calibrated rather than maximised reliance.

\subsection{Android Visibility and Permission Semantics}

Android permissions are part of a layered security model. Manifest declarations, runtime grants, AppOps state, operating-system services, and OEM behaviour do not expose one universal event stream to an ordinary third-party application. Android's permission guidance emphasises minimising permission requests and using platform-supported workflows \cite{androidpermissions2026}. Therefore, a credible review tool must state its acquisition authority and must not equate an empty result with absence of processing.

Static manifest analysis and dynamic monitoring answer different questions. Static tools can identify declared capabilities and embedded components; dynamic or taint-based systems can observe selected flows but may require instrumentation, system modification, a VPN, or elevated authority. PiOS illustrates the value and complexity of static data-flow analysis \cite{49}. Privacy Lens deliberately occupies a lighter position: it evaluates bounded audit summaries and communicates their coverage rather than claiming universal enforcement or surveillance.

\subsubsection{A Taxonomy of Mobile Privacy Evidence}

Existing approaches can be organised by the claim their evidence supports. Manifest analysis answers what capabilities an application declares and which components are packaged. It is reproducible and can be performed without executing the application, but it cannot establish that a declared capability was exercised. Static data-flow analysis attempts to connect sources and sinks in program code. It can reveal possible flows beyond manifest declarations, although reflection, native code, dynamic loading, and environment-specific paths complicate coverage. PiOS and related systems demonstrate the analytical value of tracing potential information flows while also illustrating that such analysis is specialised rather than a routine end-user feature \cite{49}.

Runtime instrumentation answers which monitored operations occurred during a particular execution. Its result depends on exercised paths, instrumentation placement, operating-system authority, and the definition of an event. Network interception answers which traffic crossed an observed interface and may expose destinations or payload structure. It cannot see purely local processing and may be limited by encryption, certificate pinning, protocol choice, and traffic that bypasses the observation path. Platform privacy dashboards can offer authoritative information for selected permissions, but their presentation and export capabilities are determined by the operating system.

Enforcement systems answer yet another question: whether a protected operation was blocked or modified under a policy. Enforcement is stronger than observation for the operation under control, but it can require device-owner status, system integration, rooting, a custom operating system, or application rewriting. These deployment choices change the threat model. A VPN monitor, for example, gains visibility by becoming an intermediary and must itself be trusted. A rooted monitor can inspect more state but weakens assumptions that apply to ordinary consumer devices.

Privacy Lens does not attempt to dominate this taxonomy. It accepts a structured summary from a bounded source and focuses on interpretation, traceability, and communication. This makes it complementary to manifest scanners, platform dashboards, enterprise management systems, and forensic monitors. A future authorised deployment could supply richer evidence through the same decision contract, but the application should continue to display the authority and coverage of that source.

\subsubsection{Historical Operations and the Ordinary-App Boundary}

The existence of an Android API is not evidence that any application can use it for any package. Access control can depend on signature permissions, usage-access grants, AppOps policy, device-owner status, process identity, Android release, and OEM modification. Historical-operation methods may be public in an SDK while useful cross-package results remain restricted in ordinary deployment. The physical-device evidence collected for this thesis exposed application-owned modes and scheduling state; it did not demonstrate unrestricted enumeration of other packages.

This boundary affects evaluation design. A complete system has an acquisition layer and a decision layer. The acquisition layer determines what the operating system or an authorised integration provides. The decision layer validates, aggregates, maps, stores, and presents that evidence. Synthetic injection can test the second layer exhaustively without proving the first. Conversely, an acquired record can demonstrate that a bridge works without proving the correctness of every decision boundary. The thesis reports these layers separately.

\subsection{Human Factors and Calibrated Trust}

Privacy notices impose substantial reading and comprehension costs \cite{12,13}. Adding another dense dashboard can reproduce the problem it intends to solve. Trustworthy warning design should therefore support quick orientation while preserving access to rationale and provenance. ISO 9241-11 frames usability through effectiveness, efficiency, and satisfaction in a specified context of use \cite{iso924111}. WCAG 2.2 and Android accessibility guidance further motivate non-colour status cues, meaningful labels, and sufficiently large touch targets \cite{wcag22,androidaccess}.

These sources support four design choices. First, the overview leads with the current evidence boundary. Second, colour is paired with iconography and text. Third, details are progressively disclosed through finding cards. Fourth, destructive actions such as clearing local findings are separated from routine navigation. These are conformance-oriented design decisions, not proof of usability; participant studies remain necessary.

\subsubsection{Cognitive Load and Progressive Disclosure}

Cognitive Load Theory distinguishes load inherent to a task from load introduced by its presentation \cite{27}. Privacy review has substantial intrinsic complexity because it combines technical events, purposes, law, expectations, and uncertainty. A product cannot remove that complexity, but it can avoid adding unnecessary navigation, decoration, or jargon. The earlier prototype mixed health dashboards, research projects, experiment controls, scoring views, and privacy findings. That structure made the user's immediate task unclear. The three-destination Revision 10 architecture reduces extraneous choice and uses a stable pattern for each finding.

Progressive disclosure is not equivalent to hiding caveats. The most important caveat---the limited reach of evidence---appears on the overview before the user initiates a review. Detail cards then reveal the observed count, source, pack, reference, rationale, and missing context. This ordering allows rapid orientation while preserving auditability. A legal citation without an explanation would be too terse; a full reproduction of statutory text on every card would be too heavy. The interface therefore supplies a short reference and offers the official source separately.

Dual-process accounts distinguish fast, heuristic judgement from slower deliberation \cite{22}. Severity colour can produce a rapid response, which is useful for prioritisation but dangerous if colour appears to communicate legal certainty. Textual labels, source chips, explicit caveats, and neutral action language are intended to interrupt an overly simple red-equals-illegal inference. This intention requires empirical testing. The thesis does not infer comprehension from the presence of these elements.

\subsubsection{Trust Calibration and Explanation Quality}

Trust is calibrated when reliance corresponds to actual capability. A system that understates uncertainty can induce over-reliance; a system that only displays caveats may become unusable and be ignored. Lee and See describe the importance of appropriate reliance on automation \cite{34}. In Privacy Lens, calibrated trust is pursued through claim symmetry: a positive-looking result and a warning must both show their evidence basis. A no-technical-concern state means that no rule fired within the admitted evidence, not that processing was lawful. A needs-review state means that a rule fired, not that a violation occurred.

Explanation quality has at least four dimensions. Fidelity asks whether the explanation corresponds to the implemented rule. Completeness asks whether the important inputs and limits are represented. Comprehensibility asks whether the intended audience can understand it. Actionability asks whether the explanation supports a reasonable next step. The deterministic rule and stored pack metadata support fidelity. Missing-context questions support completeness. The card hierarchy aims at comprehensibility. Suggested review questions support actionability. Only the first two dimensions are strongly evaluated in the present work; comprehension and actionability need participant evidence.

\subsubsection{Accessibility as an Evidence Requirement}

Accessibility affects whether transparency is available in practice. WCAG 2.2 requires that information is not conveyed through colour alone and includes criteria concerning contrast and target size \cite{wcag22}. Android guidance recommends meaningful content descriptions and interaction targets that are large enough to operate reliably \cite{androidaccess}. The implementation pairs icons and words with colour and supplies accessibility labels for primary controls.

These measures are design evidence, not conformance certification. True accessibility evaluation should include screen-reader order, announcements after asynchronous review, large-font reflow, switch access, keyboard focus, high-contrast settings, reduced motion, colour-vision deficiencies, and comprehension with users who rely on assistive technology. The distinction mirrors the wider thesis: source inspection can establish that a label exists, whereas a user study is needed to establish that the resulting interaction is effective.

\subsection{Replaceable Rule Systems}

Hard-coding statutory names into a decision engine creates maintenance and validity risks. A legal text may change, another jurisdiction may use different concepts, and a translated label may not preserve the same meaning. The appropriate software pattern is a stable engine contract with versioned, independently reviewable rule data. A pack should declare its jurisdiction, version, official source, caveat, supported signals, thresholds, references, and explanatory text. The engine should reject unknown packs and preserve the selected pack identifier in every finding.

Replaceability does not mean equivalence. Two packs cannot be treated as interchangeable simply because they implement the same TypeScript interface. Each requires independent legal review, localisation, validation fixtures, and governance over updates. The included ``Research Baseline'' pack demonstrates switching without representing a second jurisdictional legal interpretation.

\subsubsection{Pack Lifecycle and Review Governance}

A credible pack lifecycle begins before implementation. The maintainer identifies the intended jurisdiction, material and territorial scope, regulated actors, effective date, official consolidated source, language versions, and the technical signals actually available. A qualified reviewer then maps signals to questions, not verdicts, and records assumptions and exclusions. Engineering converts the reviewed mapping into immutable data and fixtures. Localisation is reviewed separately because legally important distinctions can be lost even when the code remains valid.

Release governance should treat a pack as a versioned dependency. A substantive legal amendment, authoritative interpretation, corrected translation, changed threshold, or altered missing-context question may require a new pack version. Findings retain the version used at evaluation time. Updating the active pack should not rewrite historical records. If a migration is needed, it should create a new evaluation linked to the old one rather than silently changing the original rationale.

Pack selection also raises a scope question. Device location alone is not a reliable selector because applicable law can depend on establishment, targeting, residence, processing context, contracts, and sector. Revision 10 therefore requires explicit user selection from registered packs and presents the jurisdiction label and caveat. Automatic selection could be added only with a defensible scope model and an override.

\subsubsection{Rule-Based and Learned Decision Models}

Rule-based systems are attractive for a legal-support interface because a reviewer can reconstruct the condition that fired. A fixed threshold also has obvious limitations: it can be bypassed by splitting activity across windows, may not represent purpose, and may treat different application categories alike. Learned anomaly detection can adapt to patterns, but it introduces training-data bias, distribution shift, opacity, and uncertainty. A hybrid system could combine transparent guardrails with contextual or personalised models, provided that model provenance and uncertainty remain visible.

The present engine intentionally uses explicit research thresholds. The contribution lies in safe admission, temporal isolation, versioning, explanation, and boundary disclosure rather than in claiming optimal thresholds. Population calibration would require representative applications, ground-truth definitions, stratification by purpose and category, independent annotation, and validation on data not used to choose parameters. Without that evidence, a mathematically sophisticated adaptive threshold could be less trustworthy than a simple rule whose limits are stated.

\subsubsection{Automated Validation and Oracle Independence}

Randomised tests can explore more input combinations than a short hand-written suite, but their evidential value depends on the generator and oracle. Uniform random counts may seldom exercise exact thresholds. A generator that only creates well-typed values cannot test the runtime boundary. An oracle copied from the production helper may reproduce the same defect. The evaluation therefore combines fixed seeds, boundary values, malformed types, permission mutation, temporal constructions, and an explicit expected-result calculation.

Fuzzing and property-oriented testing are especially useful for invariants. Counts should remain finite and non-negative. Results should not depend on the order of unrelated package histories. Replaying an identical window should not increase totals. Overlapping aggregate windows should not be added. A simulator source should never become a native-observation source. Unknown packs should not load. These properties express safety behaviour more clearly than a single accuracy number.

Statistical language must also match the sample. A deterministic 1,000-case run provides strong regression coverage for its generator, but it is not a random sample of real applications or alleged infringements. Confidence intervals applied to such a run would describe the generated distribution, not the mobile ecosystem. Revision 10 reports the confusion matrix descriptively and reserves population claims for future independent datasets.

\subsection{Store Delivery as a Trust Boundary}

Application-store readiness includes more than a successful APK. Google Play uses Android App Bundles for new application delivery, requires current target API levels, and requires developers to complete a Data safety declaration \cite{androidbundle2026,playtarget2026,playdatasafety2026}. Signing identity, a public privacy-policy URL, support contact, screenshots, listing text, and console review remain owner-controlled steps. A research build can be an initial submission candidate while still being unpublishable until these external obligations are completed.

\subsection{Synthesis}

The literature yields four gaps addressed by the prototype: permission signals are easily over-interpreted; acquisition coverage is often hidden; legal mappings are commonly coupled to product code; and engineering success is frequently conflated with deployment or legal validity. The following chapter converts these gaps into explicit system requirements and evidence contracts.
